\documentclass[11pt,a4paper]{article}
\usepackage[utf8]{inputenc}
\usepackage[T1]{fontenc}
\usepackage{amsmath,amssymb,amsthm}
\usepackage{listings}
\usepackage{geometry}
\usepackage{hyperref}
\geometry{margin=2.5cm}

\newtheorem{theorem}{Theorem}
\newtheorem{lemma}[theorem]{Lemma}

\lstset{
  language=C++,
  basicstyle=\ttfamily\small,
  keywordstyle=\bfseries,
  breaklines=true,
  numbers=left,
  numberstyle=\tiny,
  frame=single,
  tabsize=2
}

\title{IOI 2022 --- Insects}
\author{}
\date{}

\begin{document}
\maketitle

\section{Problem Statement}

There are $N$ insects, each of some unknown type.  A machine supports
three operations:
\begin{itemize}
  \item $\texttt{move\_inside}(i)$: place insect $i$ in the machine.
  \item $\texttt{move\_outside}(i)$: remove insect $i$ from the machine.
  \item $\texttt{press\_button}()$: return the maximum frequency among all
        types currently in the machine.
\end{itemize}
Determine the cardinality of the rarest insect type.

\section{Solution}

\subsection{Step 1: Count distinct types}

Insert insects one by one.  After each insertion, if
$\texttt{press\_button}() > 1$, remove the insect (it is a duplicate).
The number of insects remaining in the machine equals the number of
distinct types~$d$.  Then clear the machine.

This step uses at most $2N$ operations ($N$ insertions + at most $N$ removals
+ $N$ button presses + $d$ removals to clear).

\subsection{Step 2: Binary search on the answer}

Binary search on $k \in [1, \lfloor N/d \rfloor]$, testing whether the
minimum type frequency is at most $k$.

\paragraph{Check for a given $k$.}
Insert insects one by one, rejecting any insect whose insertion would
cause $\texttt{press\_button}() > k$.  After processing all $N$ insects, the
machine contains exactly $\min(k, \mathrm{count}(t))$ insects of each type~$t$.
The total inside is $T = \sum_t \min(k, \mathrm{count}(t))$.

If $T < d \cdot k$, then some type has fewer than $k$ insects, so the
minimum frequency is $\le k$.  Otherwise every type has count $\ge k$
and we set $\mathrm{lo} = k + 1$.

Each check uses $O(N)$ operations.  The binary search has
$O(\log(N/d))$ iterations.

\subsection{Total operation count}

$O(N \log(N/d))$ operations overall.  For the IOI limits
($N \le 2000$), this is at most roughly $2000 \times 11 \approx 22000$,
well within the allowed budget.

\section{Implementation}

\begin{lstlisting}
#include <bits/stdc++.h>
using namespace std;

void move_inside(int i);
void move_outside(int i);
int press_button();

int min_cardinality(int N) {
    // Step 1: count distinct types
    vector<bool> inside(N, false);
    int num_types = 0;
    for (int i = 0; i < N; i++) {
        move_inside(i);
        if (press_button() > 1) {
            move_outside(i);
        } else {
            inside[i] = true;
            num_types++;
        }
    }
    // Clear the machine
    for (int i = 0; i < N; i++)
        if (inside[i]) {
            move_outside(i);
            inside[i] = false;
        }

    // Step 2: binary search on the rarest frequency
    int lo = 1, hi = N / num_types;
    while (lo < hi) {
        int mid = (lo + hi) / 2;

        // Check: is the minimum frequency <= mid?
        int total_inside = 0;
        vector<bool> in_machine(N, false);
        for (int i = 0; i < N; i++) {
            move_inside(i);
            if (press_button() > mid) {
                move_outside(i);
            } else {
                in_machine[i] = true;
                total_inside++;
            }
        }

        bool has_rare = (total_inside < (long long)num_types * mid);

        // Clean up
        for (int i = 0; i < N; i++)
            if (in_machine[i])
                move_outside(i);

        if (has_rare)
            hi = mid;
        else
            lo = mid + 1;
    }
    return lo;
}
\end{lstlisting}

\section{Complexity}

\begin{itemize}
  \item \textbf{Operations:} $O(N \log(N/d))$ where $d$ is the number of
        distinct types.  Each binary-search iteration uses $O(N)$
        move/press calls.
  \item \textbf{Space:} $O(N)$.
\end{itemize}

\end{document}
